\SourceLeaf{536}senkrechten Spalte jeder einzelnen Tafel die drei niedrigsten Ziffern des Quadrats und im Inneren die vier oder fünf höheren Ziffern des Quadrats.
\par Die Tafel ‘\textit{Indizes der Primzahlen im höheren Zahlenbereich}’ enthält in der obersten waagerechten Reihe den jeweiligen Modul, in der zweiten Reihe die verwendete Basis, in der ersten senkrechten Spalte die Restzahlen und im Inneren die Indizes. Die Sterne $*$ bezeichnen die Nullreste.
\par Die Handschrift des Fragments ‘\textit{Achter Abschnitt. Ausführlichere Betrachtung einiger Untersuchungen zur Kreisteilung}’ scheint aus der Zeit der Umarbeitung der ‘Analysis residuorum’ in die ‘Disquisitiones arithmeticae’ zu stammen. Die Briefe von \textsc{Gauss} an \textsc{Dirichlet} bestätigen die Ansicht, dass \textsc{Gauss} auch in der höheren Arithmetik erheblich mehr entdeckt hat, als sich im Nachlass findet.
\begin{flushright}\textsc{Schering}.\end{flushright}
\BookRule{.25\textwidth}
